\documentclass[11pt,a4paper]{ctexart}
\usepackage[margin=2.2cm]{geometry}
\usepackage{amsmath,amssymb,amsthm,mathtools,booktabs,hyperref,microtype}
\hypersetup{colorlinks=true,linkcolor=black,urlcolor=blue,citecolor=black}
\newtheorem{theorem}{定理}
\newtheorem{lemma}{引理}
\newtheorem{corollary}{推论}
\renewcommand{\proofname}{证明}
\newcommand{\F}{F_{3,D}}
\newcommand{\vthree}{v_3}
\title{\bfseries 固定整数多项式对 \(F_{3,D}\) 的\\实现能力与最低次数}
\author{研究札记 · F3D-POLY}
\date{2026 年 9 月 29 日}
\begin{document}
\maketitle

\begin{abstract}
在定义域 \(\mathcal D=\{(n,d)\in\mathbb Z^2:n\ne\pm d\}\) 上，令
\[
F_{3,D}(n,d)=v_3(n+d)-v_3(n-d).
\]
本文首先给出一个四次整数多项式对，它在全定义域上精确检测 \(F_{3,D}=0\)，并由线性层平移得到任意指定层的四次检测器。随后完整分类固定整数多项式对能够对 \(F_{3,D}\) 数值实现的单输入函数：恰好是正负两端分别最终成为整数仿射函数的整数值函数。证明在和差比 \(x=(n+d)/(n-d)\) 上将问题化为有理函数三进赋值，并用有限个正部折点显式构造充分性。最后证明四个基础操作的最低次数分别为 \(1,2,3,4\)：层平移、绝对值、取正部与精确零层检测。由分类立即得到 \(t^2\)、奇偶振荡以及固定双输入 \(F\)-乘法器的不可实现性。
\end{abstract}

\section{和差比坐标}
令
\[
u=n+d,\qquad v=n-d,\qquad x=\frac uv.
\]
在 \(\mathcal D\) 中 \(u,v\ne0\)，所以
\[
\boxed{\F(n,d)=v_3(x).}
\]
任意非零有理数 \(x=a/b\) 都可由整数输入
\[
n=a+b,\qquad d=a-b
\]
实现，因此 \(F_{3,D}\) 的值域问题等价于 \(\mathbb Q^\times\) 上的三进赋值问题。

\section{四次精确层检测}
定义
\[
E_0(n,d)=
\left(
5n^4+22n^2d^2+5d^4,\,
(n^2-d^2)^2
\right).
\]
令
\[
A=u^4+u^2v^2+v^4,\qquad B=u^4+v^4.
\]
直接展开得
\[
E_0(n,d)=(A+B,A-B),
\]
故
\[
\F(E_0(n,d))=v_3(A)-v_3(B).
\]

写
\[
a=v_3(u),\quad b=v_3(v),\quad m=\min(a,b).
\]
若 \(a\ne b\)，\(A\) 与 \(B\) 的最低赋值项均唯一，故
\[
v_3(A)=v_3(B)=4m.
\]
若 \(a=b=m\)，写
\[
u=3^mu_0,\qquad v=3^mv_0,
\]
其中 \(u_0,v_0\) 为单位。此时
\[
v_3(B)=4m
\]
因为 \(u_0^4+v_0^4\equiv2\pmod3\)。

再写
\[
u_0^2=1+3r,\qquad v_0^2=1+3s.
\]
模 \(9\) 有
\[
u_0^4+u_0^2v_0^2+v_0^4
\equiv3\pmod9.
\]
所以
\[
v_3(A)=4m+1.
\]
因此
\[
\boxed{
\F(E_0(n,d))
=
\mathbf1_{\{\F(n,d)=0\}}.
}
\]
并且因
\[
A-B=u^2v^2>0,\qquad B>0,
\]
输出始终满足 \(N>D>0\)。

对任意 \(r\in\mathbb Z\)，可取显式线性整数多项式对 \(S_r\) 使
\[
\F(S_r(n,d))=\F(n,d)+r.
\]
于是
\[
E_r=E_0\circ S_{-r}
\]
仍为四次，并满足
\[
\boxed{
\F(E_r(n,d))
=
\mathbf1_{\{\F(n,d)=r\}}.
}
\]

\section{分类定理：必要性}
设固定整数多项式对 \((P,Q)\) 全域实现 \(g:\mathbb Z\to\mathbb Z\)。令
\[
A=P+Q,\qquad B=P-Q.
\]

取
\[
(n,d)=(3^m+1,3^m-1).
\]
则 \(\F(n,d)=m\)。代入 \(A,B\) 后得到关于 \(3^m\) 的两个固定非零整数多项式。

若
\[
H(X)=c_jX^j+\cdots,\qquad c_j\ne0
\]
且 \(j\) 是最低非零次数，则当 \(m\) 充分大，
\[
v_3(H(3^m))=jm+v_3(c_j),
\]
因为最低次数项最终具有唯一最小赋值。

因此
\[
g(m)=a_+m+b_+
\]
对充分大 \(m\) 成立。

再取
\[
(n,d)=(3^m+1,1-3^m),
\]
此时 \(\F(n,d)=-m\)，同理得到
\[
g(t)=a_-t+b_-
\]
对充分小 \(t\) 成立。

所以任何可实现函数在正负两端都必须最终整数仿射。

\section{分类定理：充分性}
定义
\[
\rho(x)=\frac{x(x^2+1)}{x^3-x+1}.
\]
它在 \(\mathbb Q^\times\) 上无零点或极点，并满足
\[
\boxed{
v_3(\rho(x))=\max(v_3(x),0).
}
\]
这是直接按 \(t=v_3(x)>0,=0,<0\) 三种情形比较最低赋值项得到的。

现在设 \(g\) 两端最终仿射，负端为
\[
g(t)=a_-t+b_-.
\]
定义离散二阶差分
\[
c_r=g(r+1)-2g(r)+g(r-1).
\]
仅有限多个 \(c_r\ne0\)，并有
\[
g(t)=a_-t+b_-+\sum_r c_r(t-r)_+.
\]
因此定义
\[
\boxed{
\mathcal R_g(x)=
3^{b_-}x^{a_-}
\prod_{r:c_r\ne0}\rho(3^{-r}x)^{c_r}.
}
\]
赋值加性给出
\[
v_3(\mathcal R_g(x))=g(v_3(x)).
\]

写 \(\mathcal R_g=A/B\)，清有理系数分母并齐次化：
\[
\widetilde A(u,v)=v^N A(u/v),\qquad
\widetilde B(u,v)=v^N B(u/v).
\]
再定义
\[
P=L(\widetilde A+\widetilde B),\qquad
Q=L(\widetilde A-\widetilde B)
\]
并取 \(L\) 清分母。因为所有构造因子在 \(\mathbb Q^\times\) 上非零，输出始终留在定义域，并且
\[
\F(P,Q)=g(\F(n,d)).
\]

由此得到：
\[
\boxed{
g\text{ 可由固定整数多项式对全域实现}
\iff
g\text{ 两端最终整数仿射}.
}
\]

\section{最低次数}
次数指输出两个整数多项式的最大总次数。

\subsection{层平移：1}
线性 \(S_r\) 已实现 \(t\mapsto t+r\)。常数多项式不可能产生随 \(t\) 变化的输出，所以最低次数为
\[
\boxed1.
\]

\subsection{绝对值：2}
有理函数
\[
R_{\rm abs}(x)=\frac{x}{x^2+1}
\]
满足
\[
v_3(R_{\rm abs}(x))=|v_3(x)|.
\]
齐次化给出二次实现。

若次数至多 \(1\)，正端斜率 \(+1\) 要求零点阶差为 \(1\)，负端斜率 \(-1\) 又要求分子分母次数差为 \(-1\)，在线性次数内不可能同时满足。故最低次数为
\[
\boxed2.
\]

\subsection{取正部：3}
\(\rho\) 给出三次实现。

若存在二次实现，非齐次候选可通过共同 \(3^m\) 缩放提取最低齐次部分，降为次数至多 \(2\) 的既约有理函数 \(A/B\)。正端斜率 \(1\) 强迫分子恰含一个 \(x\) 因子而分母在 \(0\) 非零；负端输出 \(0\) 强迫 \(\deg A=\deg B\)。公共次数为 \(1\) 时分母必有非零有理根；公共次数为 \(2\) 时分子除去 \(x\) 后剩一个一次因子，也必有非零有理根。两者都违反全域可容许性。因此最低次数为
\[
\boxed3.
\]

\subsection{精确零层：4}
四次构造已经给出上界。

假设存在次数至多 \(3\) 的实现。通过同样的非齐次降阶，得到次数至多 \(3\) 的既约有理函数
\[
R=A/B
\]
满足
\[
v_3(R(x))=
\begin{cases}
1,&v_3(x)=0,\\
0,&v_3(x)\ne0.
\end{cases}
\]

有理数逼近首先排除 \(A,B\) 的任何非零 \(\mathbb Q_3\) 根。正负尾都为零，又强迫：
\[
A(0),B(0)\ne0,\qquad
\deg A=\deg B,
\]
且分子分母常数项赋值相同、首项系数赋值相同。

次数 \(1\) 会直接产生 \(\mathbb Q_3\) 根，所以只剩次数 \(2\) 或 \(3\)。此时无一次因子意味着 \(A,B\) 在 \(\mathbb Q_3[x]\) 中不可约。完备赋值域绝对值向有限扩张唯一延拓，所以同一不可约多项式的共轭根具有相同赋值。韦达公式与首尾赋值相等进一步迫使 \(A,B\) 的根具有同一三进半径。

若该半径不是单位半径，则任意有理单位 \(x\) 到所有根的赋值由最小值公式固定，得到 \(v_3(A/B)=0\)，矛盾。因此所有根都是单位。归一化后
\[
A,B\in\mathbb Z_3[x]
\]
且首尾系数为单位。

单位层要求输出 \(1\)，特别地
\[
A(1)\equiv A(-1)\equiv0\pmod3.
\]
次数至多 \(3\) 时，两个不同根不可能同时都是重根，因此至少有一个简单根。Hensel 引理把它提升成非零 \(\mathbb Q_3\) 根，与前面的根排除矛盾。

所以精确检测的最低次数是
\[
\boxed4.
\]

\section{不可实现性}
分类定理立即排除
\[
t^2,\qquad (-1)^t,\qquad \left\lfloor\frac t2\right\rfloor,
\]
因为它们至少一端不最终整数仿射。

同样，不存在固定双输入整数多项式构造恒实现
\[
F(\Phi(X,Y))=F(X)F(Y).
\]
否则令 \(Y=X\)，便得到单输入 \(t^2\)，矛盾。

\section{状态与边界}
本文分类的是固定整数系数多项式对的全域实现能力。允许条件分支、循环、输入相关公式或只在有限子域成立时，不属于本定理范围。

本稿状态为 \texttt{PAPER-AUDITED}，尚未完成 Lean 形式化。

\begin{thebibliography}{9}
\bibitem{mathlib-val}
Mathlib, \emph{Mathlib.NumberTheory.Padics.PadicVal.Basic}, current online documentation.

\bibitem{mathlib-hensel}
Mathlib, \emph{Mathlib.NumberTheory.Padics.Hensel}, current online documentation.

\bibitem{MIT}
Andrew V. Sutherland, \emph{18.785 Number Theory I, Lecture 10: Extensions of complete DVRs}, MIT, 2021. Theorem 10.4 and Remark 10.5.
\end{thebibliography}

\end{document}
